\section{如何学好刑诉法}
\begin{enumerate}
  \item 7 月份进入了法考备战的黄金时间，需要对重点的考点进行专项的背诵记忆。
  \item 客观题对考点的归纳总结背诵，对选项的判断能力。
  \item 主观题考的是对案情的阅读能力，对考点的识别能力，法条搜索能力，结论的书写能力，论证说理的能力。
  \item 案例化，法条化，实务操作能力，关注身边的真实的法律新闻和案例。
  \item 新法必考，新增必考，重者恒重，热点必考（热点案例/改革热点）
  \item 上诉不加刑原则，管辖制度，证据制度，辩护制度，认罪认罚从宽制度，附条件不起诉制度。
  \item 重要的法律文件。
        \begin{enumerate}
          \item 刑诉法。
          \item 刑诉解释、高检规则、公安部规定。
          \item 专项问题找专规。排非规定、认罪认罚指导意见。
          \item 2023 年。
                \begin{enumerate}
                  \item 未成年人犯罪记录封存的实施办法。
                  \item 取保候审的若干问题的规定。
                  \item 《关于办理信息网络犯罪案件适用刑事诉讼程序若干问题的意见》。
                        \begin{enumerate}
                          \item 管辖。沾边就管。
                          \item 网络犯罪案件。可以先并后分。公安可以并案管辖，侦查完毕后分开移送给检察院。
                          \item 侮辱、诽谤取证可以要求公安提供协助 \yjt 手续合法 \yjt 应当协助。\al{A 市公安请求 B 市公安协助取证，B 市公安应当协助取证 \cuo}。
                          \item 雷同证据可按比例收集。
                          \item 不科学收集的，应当由原取证机关作出补正或重新取证。
                          \item 不能排除合理怀疑，做有利于被告的推定。
                          \item 无法核实每一笔赃款，可视为卡里有多少算多少，被告人能自证的除外。
                        \end{enumerate}
                \end{enumerate}
        \end{enumerate}
  \item 改革热点。
        \begin{enumerate}
          \item 以审判为中心的诉讼制度改革。
          \item 控辩式的审判模式改革。
          \item 认罪认罚从宽制度改革。
          \item 缺席审判改革。
        \end{enumerate}
  \item 全局框架。三段。
        \begin{enumerate}
          \item 10 项原理和 10 项制度。基本原则、专门机关和诉讼参加人、管辖、回避、辩护和代理、证据、强制措施、附带民事、期间送达。
          \item 立案、侦查、起诉、审判（一审、二审、复核、再审）、执行。
          \item 5 个特别程序加 1 个涉外程序。未成年人、公诉和解、缺席审判、没收、强制医疗、涉外刑诉。
        \end{enumerate}
\end{enumerate}
